\section{Passage path: can it move through both chambers and then work?}
\label{sec:passage}

The eighth question closes the loop. Even a good bill needs a path, and a member
wants to know the bill can both pass and then be implemented without embarrassment.

[DRAFTING INSTRUCTIONS: In the full stage, lay out the path from
\texttt{a9-research-matrix.tex} and the bill's effective-date and rulemaking
provisions in \texttt{s3-amendment.tex}: introduction and referral, subcommittee and
full-committee markup, House floor passage, a Senate companion and the same
sequence, amendment exchange or conference, and enrollment; the vote thresholds (a
House majority, Senate cloture at 60, Senate passage at 51); and the implementation
runway (final regulations within 365 days, effective the first quarter at least one
year after enactment, 180 days for enrolling investigations). Close by showing the
platform is ready to implement what the bill codifies
\cite{Kawchak2026NationalPlatformPhysicalAI, Kawchak2026FullyAutomatedSponsorPhysical}.
Reference the legislative process \cite{congress2024lawsmade}. This section carries
the legislative-journey figure (diagram 06), the Senate-companion sequence (diagram
15), the markup gitGraph (diagram 16), the vote-thresholds requirement diagram
(diagram 17), and the capstone (diagram 20).]

\begin{psgfig}
\node[psgone] (h) at (0,0) {House\\passage};
\node[psgevid] (s) at (4.0,0) {Senate\\passage};
\node[psgact] (l) at (8.0,0) {Enacted\\section 515D};
\draw[psgedge] (h)--(s); \draw[psgedge] (s)--(l);
\end{psgfig}
\psgcaption{Figure (placeholder). The path from two chambers to one statute. [FULL
STAGE: replace with diagram 20 capstone and the synthesis table.]}
